% -------------------------
% Schriftart und Layout
% -------------------------
\usepackage[printonlyused]{acronym}
\usepackage{times}          % Verwende die Times-Schriftart für das gesamte Dokument
\usepackage{setspace}       % Erlaubt die Anpassung des Zeilenabstands im Dokument
\usepackage[utf8]{inputenc} % UTF-8 Zeichencodierung
\usepackage[T1]{fontenc}    % Korrekte Schriftkodierung für Umlaute
\usepackage[sorting=none, backend=biber]{biblatex}  % Backend auf biber setzen
\usepackage{pdfpages}

% -------------------------
% Grafische und Layout-Pakete
% -------------------------

\usepackage{tikz}           % Ermöglicht das Erstellen von Diagrammen, Zeichnungen und Illustrationen
\usepackage{graphicx}       % Zum Einfügen und Skalieren von Bildern im Dokument
\usepackage{geometry}       % Anpassung der Seitenränder und Layout des Dokuments
\usepackage[german]{babel}  % Deutsche Sprachunterstützung für Trennungen, Datum etc.
\sloppy
\usepackage{datetime}       % Erlaubt die Formatierung von Datumsangaben im Dokument
\usepackage[hidelinks]{hyperref} % Erzeugt klickbare Links im Dokument, ohne farbige Umrahmungen
\usepackage[absolute]{textpos} % Ermöglicht das absolut positionierte Platzieren von Text auf der Seite
\usepackage{extra/xurl}     % Bessere URL-Trennung
\usepackage{dirtree}        % Verzeichnisstruktur

% -------------------------
% Kopf- und Fußzeilenverwaltung
% -------------------------

\usepackage{fancyhdr}       % Ermöglicht die Anpassung von Kopf- und Fußzeilen
\setlength{\headheight}{15pt} % Stellt sicher, dass genug Platz für die Kopfzeile vorhanden ist

% -------------------------
% Tabellen und Listen
% -------------------------

\usepackage{array}          % Bietet erweiterte Möglichkeiten zur Tabellenformatierung
\usepackage{caption}        % Ermöglicht die Anpassung von Bild- und Tabellenunterschriften
\usepackage{float}          % Präzisere Kontrolle der Platzierung von Gleitobjekten (Tabellen, Abbildungen)
\usepackage{supertabular}   % Erstellung von sehr langen Tabellen, die über mehrere Seiten gehen
\usepackage{tabularx}       % Für automatische Anpassung der Spaltenbreite
\usepackage{makecell}

% -------------------------
% Quellcode-Darstellung
% -------------------------
\usepackage{listings}       % Darstellung von Quellcode mit Syntax-Highlighting
\usepackage{xcolor}         % Ermöglicht die Verwendung von Farben im Dokument

% -------------------------
% Referenzierung und Zitate
% -------------------------

\usepackage{cleveref}       % Automatisierte Verweise mit intelligenten Namen, z.B. "Abschnitt" statt "Section"
\usepackage{csquotes}       % Bessere Handhabung und Formatierung von Zitaten

% -------------------------
% Sektionstitel anpassen
% -------------------------

%\usepackage{titlesec}       % Erlaubt eine detaillierte Anpassung der Kapitel- und Abschnittstitel

% -------------------------
% Nummerierungsanpassungen
% -------------------------

\usepackage{chngcntr}       % Ermöglicht das Anpassen der Nummerierung von Tabellen, Abbildungen etc.

% -------------------------
% Definition der Seitenränder
% -------------------------

% Definiert die Seitenränder des Dokuments als Makros, damit sie leicht angepasst werden können
\newcommand{\leftMargin}{3cm}
\newcommand{\rightMargin}{2.5cm}
\newcommand{\topMargin}{2.5cm}
\newcommand{\bottomMargin}{2cm}

% -------------------------
% Pfad zu Bildern
% -------------------------

\graphicspath{ {images/} }  % Setzt den Standardpfad für Bilddateien im Dokument

% -------------------------
% Anwendung der Seiteneinstellungen
% -------------------------

\geometry{
  a4paper,               % Papiergröße: A4
  left=\leftMargin,      % Linker Seitenrand
  right=\rightMargin,    % Rechter Seitenrand
  top=\topMargin,        % Oberer Seitenrand
  bottom=\bottomMargin,  % Unterer Seitenrand
  includehead            % Schließt die Kopfzeile in den oberen Rand ein
}

% -------------------------
% Zeilenabstand
% -------------------------

\setstretch{1.2}            % Setzt den Zeilenabstand auf 1,2-fach für das gesamte Dokument

% -------------------------
% Kopfzeileneinstellungen
% -------------------------

\fancyhead{}                            % Leert alle vordefinierten Kopfzeilen-Einstellungen
\fancyhead[L]{\nouppercase{\leftmark}}  % Setzt die linke Kopfzeile ohne automatische Großschreibung des Textes
\renewcommand{\headrulewidth}{0.4pt}    % Fügt einen 0,4pt dicken Strich unter der Kopfzeile hinzu

% -------------------------
% Seitenstil für Seiten ohne Kopfzeile
% -------------------------

\fancypagestyle{noheader}{
  \fancyhead{}              % Keine Kopfzeile
  \renewcommand{\headrulewidth}{0pt} % Kein Strich unter der Kopfzeile
}

% -------------------------
% Anpassung des Abschnittstitels
% -------------------------

%\titleformat{\section}[block]
%  {\normalfont\Large\bfseries}{\thesection}{1em}{}
%% Setzt das Format für Abschnittstitel: normale Schrift, große, fettgedruckte Titel, nummeriert
%
%% Abstand vor und nach dem Abschnittstitel
%\titlespacing*{\section}{0pt}{\baselineskip}{\baselineskip}
%% Definiert den Abstand vor und nach dem Abschnittstitel

% -------------------------
% Spezifischer Seitenumbruch
% -------------------------

% Definiert einen neuen Befehl für einen Seitenumbruch mit einem spezifischen Seitenstil (ohne Kopfzeile)
\newcommand{\sectionbreak}{\clearpage\thispagestyle{noheader}}

% -------------------------
% Nummerierung der Abbildungen, Tabellen und Listings nach Abschnitten
% -------------------------

\counterwithin{figure}{section}  % Abbildungen werden nach den Abschnitten nummeriert
\counterwithin{table}{section}   % Tabellen werden nach den Abschnitten nummeriert
\AtBeginDocument{
  \counterwithin{lstlisting}{section} % Quellcode wird nach den Abschnitten nummeriert
  \crefname{lstlisting}{Quellcode}{Quellcodes} % Verweise auf Quellcodes richtig benennen
  \Crefname{lstlisting}{Quellcode}{Quellcodes}
}

% -------------------------
% Umbenennung der Bezeichnungen für Quellcode
% -------------------------

\renewcommand{\lstlistingname}{Quellcode}        % Standardbezeichnung für Quellcode
\renewcommand{\lstlistlistingname}{Quellcodeverzeichnis} % Standardbezeichnung für das Quellcodeverzeichnis

% -------------------------
% Definitionen für clevere Referenzen
% -------------------------

\crefname{lstlisting}{Quellcode}{Quellcodes}  % Intelligente Referenzierung für Quellcodes
\crefname{section}{Abschnitt}{Abschnitte}     % Intelligente Referenzierung für Abschnitte
\crefname{figure}{Abbildung}{Abbildungen}     % Intelligente Referenzierung für Abbildungen
\crefname{table}{Tabelle}{Tabellen}           % Intelligente Referenzierung für Tabellen
\crefname{equation}{Gleichung}{Gleichungen}   % Intelligente Referenzierung für Gleichungen

% -------------------------
% Setup für Akronyme
% -------------------------

%\acsetup{make-links} % Verlinkt die Definitionen von Akronymen automatisch im Text

% -------------------------
% Anpassung der Farben für den Quellcode
% -------------------------

\definecolor{keywordcolor}{rgb}{0.13, 0.13, 1.0}   % Definiert die Farbe für Schlüsselwörter im Quellcode
\definecolor{commentcolor}{rgb}{0.0, 0.5, 0.0}     % Definiert die Farbe für Kommentare im Quellcode
\definecolor{stringcolor}{rgb}{0.58, 0.0, 0.82}    % Definiert die Farbe für Zeichenketten im Quellcode
\definecolor{backgroundcolor}{rgb}{0.95, 0.95, 0.95} % Definiert die Hintergrundfarbe des Quellcodes

% Einstellungen für die Darstellung von Quellcode
\lstset{
    basicstyle=\ttfamily\scriptsize,   % Grundstil: Schreibmaschinenschrift, klein
    keywordstyle=\color{keywordcolor}, % Stil für Schlüsselwörter
    commentstyle=\color{commentcolor}, % Stil für Kommentare
    stringstyle=\color{stringcolor},   % Stil für Zeichenketten
    backgroundcolor=\color{backgroundcolor}, % Hintergrundfarbe für Quellcode
    numberstyle=\tiny\color{black},    % Stil der Zeilennummern: klein und schwarz
    showstringspaces=false,            % Leerzeichen in Zeichenketten nicht anzeigen
    numbers=left,                      % Zeilennummern links anzeigen
    numbersep=5pt,                     % Abstand der Zeilennummern vom Quellcode
    lineskip=-0.1pt,                   % Zeilenabstand leicht reduzieren
    breaklines=true,                   % Zeilen automatisch umbrechen, falls sie zu lang sind
    captionpos=b,                      % Bildunterschriften unter dem Quellcode platzieren
    literate=                          % Zeichenanpassungen für deutsche Umlaute und das scharfe S
        {ä}{{\"a}}1
        {ö}{{\"o}}1
        {ü}{{\"u}}1
        {ß}{{\ss}}1
        {Ä}{{\"A}}1
        {Ö}{{\"O}}1
        {Ü}{{\"U}}1
}